% =========================================================================
\chapter{Operations Narrative}
\label{app:ops}
% =========================================================================

This appendix records what shaped the campaigns without being evidence for
any claim: the infrastructure incidents and the standing rules they left
behind, and the cost mechanisms, which are engineering and are kept out of
the body so that they are not mistaken for results of the graph. Nothing
here may be cited as a result.

Two graph campaigns were flown before the graph layer was correct. They
are excluded from this thesis in full by the whitelist ruling
(\S\ref{sec:scope}): the defects they surfaced were repaired and pinned by
tests before the first admissible graph campaign, and nothing measured on
them appears in this appendix or anywhere else in this document.

% -------------------------------------------------------------------------
\section*{B.1\quad Infrastructure incidents and the rules they left behind}
% -------------------------------------------------------------------------

\begin{itemize}
  \item \textbf{Gateway routing drift.} Model names are routing keys that
        move: the same name reaching different backends, mappings changing
        underneath a campaign, and a route once judged dead coming back to
        life. \emph{Standing rule:} probe by workload shape before a
        campaign, never by a single trial call.
  \item \textbf{Gateway outage.} On the third seed the gateway
        disconnected for twelve hours mid-campaign. Affected rounds are
        identifiable because every fresh evaluation in them cost zero;
        they were isolated, and the campaign's history, audit, curve and
        journal files were washed back to a clean prefix and re-flown from
        it. Both re-flights are disclosed in the body
        (\S\ref{sec:replication}).
  \item \textbf{Machine restart.} On the second seed an unplanned restart
        killed one evolve stage and one task batch; the batch was re-flown
        from a byte-identical round-head configuration and the affected
        pair is read as a same-configuration pair (\S\ref{sec:replication}).
  \item \textbf{Format drift on the meta tier}, which breaks parsing
        silently rather than loudly.
  \item \textbf{Virtual-environment shim processes} that appear as a second
        interpreter with identical arguments. \emph{Standing rule:} check
        the parent process before judging a double launch --- killing the
        shim kills the run.
  \item \textbf{Launch spikes} producing phantom runs that must be excluded
        by inspection before scoring (registry error 14).
  \item \textbf{The recorder's crash family.} Binary content reaching the
        raw JSONL killed the evolve stage repeatedly on the graph arms; the
        counts are entered in the body as the graph layer's one measured
        operational liability (\S\ref{sec:cost}) \F{56}.
\end{itemize}

% -------------------------------------------------------------------------
\section*{B.2\quad Cost mechanisms: harness-level engineering}
\label{app:cost}
% -------------------------------------------------------------------------

Four mechanisms reduced the price of a campaign. None of them reads the
execution graph; each runs on the unmodified official system, and
\S\ref{sec:cost} declines to credit them to GHX for that reason. They are
recorded here because the replication of \S\ref{sec:replication} would not
have been affordable without them.

\paragraph{No-op round batching.} A round in which nothing shipped re-runs
a configuration that has already been measured. Under the official loop
the whole bed is re-run anyway. Restricting a no-op round to a 25-task
audit batch --- the tasks that flipped most recently, chosen by a fixed
rule --- and carrying the rest forward cut the per-round task cost from
\$36.3 to \$11.8 on the whitelist graph campaign \F{13}. The design
consequence is disclosed in \S\ref{sec:gain-face}: a delta that straddles
an audit window is a partial re-draw, so every gain-face reading restricts
to full batches.

\paragraph{Single-shot digestion.} The official Digester is an agentic
session per task. Rewriting it as one model call per task --- the same
template, the trajectory supplied whole --- removes the session overhead
entirely. Its cost effect was measured only on the excluded campaigns and
is therefore not reported; its argument-quality effects on the admissible
campaigns are reported in \S\ref{sec:what-graph-did}. The Evolver is not
digester-shaped and the same rewrite does not apply to it: its context is
a compaction sawtooth in which much of each step's input is re-paid, so
its lever is not fewer calls but not paying twice, by keeping evidence
outside the context and fetching it on demand.

\paragraph{Zero cache hits.} Across the whitelist graph campaign's 2{,}760
meta-role calls and the solver calls beside them, the provider's
\texttt{cache\_read} and \texttt{cache\_write} token counters were never
non-zero \F{56}. The provider prices a cached read at roughly one fiftieth
of a fresh one, so every input token was billed at full price. The cause
sits at the gateway and was not resolved within the project; it is
recorded as the largest cost lever with zero behavioural risk, and as a
debit against the operator rather than against either arm, since both
arms paid it equally.

\paragraph{Retention policy.} Message-retention accounting identified
context that could be released without loss --- tool results older than a
few steps, truncated to a stub unless the assistant later cites them ---
and this too is portable and was not the graph's. Its size was measured
only on the excluded campaigns and is not reported.
